\documentclass{article}
\usepackage{graphicx} % Required for inserting images
\usepackage{amsmath}
\usepackage{geometry} % Required for customizing margins

\geometry{margin=1in} % Sets 1-inch margins on all sides

\title{EMCH 360 Assignment \#1}
\author{Jacob C. Vaught}
\date{September 7th 2023}


\begin{document}

\maketitle

\section*{Problem \#1}
The force exerted by the wind on a building is given by the following equation:
\[
F = \frac{{C_D \cdot \rho \cdot V^2 \cdot A}}{2}
\]
where
\begin{itemize}
\item \( F \): force exerted by the wind, with dimensions \([M \cdot L \cdot T^{-2}]\),
\item \( \rho \): density of air, with dimensions \([M \cdot L^{-3}]\),
\item \( V \): wind speed, with dimensions \([L \cdot T^{-1}]\),
\item \( A \): cross-sectional area, with dimensions \([L^2]\),
\item \( C_D \): drag coefficient.
\end{itemize}
Substituting the dimensions into the given equation, we obtain:
\[
[M \cdot L \cdot T^{-2}] = \left( [C_D] \cdot [M \cdot L^{-3}] \cdot [L^2 \cdot T^{-2}] \cdot [L^2] \right)/2
\]
Simplifying, we find:
\[
[M \cdot L \cdot T^{-2}] = [C_D] \cdot [M \cdot L \cdot T^{-2}]
\]
Since the units on both sides of the equation cancel out, we conclude that the drag coefficient \( C_D \) is dimensionless:
\[
[C_D] = 1
\]
Thus, the dimensions of the drag coefficient are dimensionless.

\newpage % Starts a new page

\section*{Problem \#2}
Express the following quantities in SI Units.

\begin{itemize}
    \item \( a \): speed, given as \(10.2 \, \text{in./min}\)
    \item \( b \): mass, given as \(4.81 \, \text{slugs}\)
    \item \( c \): force, given as \(3.02 \, \text{lb}_f\)
    \item \( d \): acceleration, given as \(73.1 \, \text{ft/s}^2\)
    \item \( e \): force per area per time, given as \(0.0234 \, \text{lb}_f \cdot \text{s/ft}^2\)
\end{itemize}

\subsection*{Solution}

To convert these quantities into SI Units, the following conversion factors will be used:

\begin{itemize}
    \item \( 1 \, \text{inch} = 0.0254 \, \text{m} \)
    \item \( 1 \, \text{slug} = 14.5939 \, \text{kg} \)
    \item \( 1 \, \text{lb}_f = 4.44822 \, \text{N} \)
    \item \( 1 \, \text{ft} = 0.3048 \, \text{m} \)
    \item \( 1 \, \text{s} = 1 \, \text{s} \) (No conversion needed)
\end{itemize}

\subsubsection*{Part a}

\[
10.2 \, \text{in./min} = 10.2 \times \frac{0.0254 \, \text{m}}{1 \, \text{in.}} \times \frac{1 \, \text{min}}{60 \, \text{s}} = 0.00429333 \, \text{m/s}
\]

\subsubsection*{Part b}

\[
4.81 \, \text{slugs} = 4.81 \times 14.5939 \, \text{kg/slug} = 70.2238 \, \text{kg}
\]

\subsubsection*{Part c}

\[
3.02 \, \text{lb}_f = 3.02 \times 4.44822 \, \text{N/lb}_f = 13.4417 \, \text{N}
\]

\subsubsection*{Part d}

\[
73.1 \, \text{ft/s}^2 = 73.1 \times \frac{0.3048 \, \text{m}}{1 \, \text{ft}} = 22.29728 \, \text{m/s}^2
\]

\subsubsection*{Part e}

\[
0.0234 \, \text{lb}_f \cdot \text{s/ft}^2 = 0.0234 \times \frac{4.44822 \, \text{N}}{1 \, \text{lb}_f} \times \frac{1}{(0.3048 \, \text{m/ft})^2} = 1.1204 \, \text{N s/m}^2
\]
\newpage % Starts a new page

\section*{Problem \#3}
Express the following quantities in BG Units.

\begin{itemize}
    \item \( a \): distance, given as \(14.2 \, \text{km}\)
    \item \( b \): pressure, given as \(8.14 \, \text{N/m}^3\)
    \item \( c \): density, given as \(1.61 \, \text{kg/m}^3\)
    \item \( d \): viscosity, given as \(0.0320 \, \text{N.m/s}\)
    \item \( e \): velocity, given as \(5.67 \, \text{mm/hr}\)
\end{itemize}

\subsection*{Solution}

To convert these quantities into BG Units, the following conversion factors will be used:

\begin{itemize}
    \item \( 1 \, \text{km} = 0.621371 \, \text{mile} \)
    \item \( 1 \, \text{N/m}^3 = 157.087 \, \text{lb}_f/\text{ft}^3 \)
    \item \( 1 \, \text{kg/m}^3 = 0.062428 \, \text{slug/ft}^3 \)
    \item \( 1 \, \text{N.m/s} = 723.569 \, \text{lb}_f.ft/\text{s} \)
    \item \( 1 \, \text{mm/hr} = 0.00328084 \, \text{in/s} \)
\end{itemize}

\subsubsection*{Part a}

\[
14.2 \, \text{km} = 14.2 \times 0.621371 \, \text{mile/km} = 8.82347 \, \text{mile}
\]

\subsubsection*{Part b}

\[
8.14 \, \text{N/m}^3 = 8.14 \times 157.087 \, \text{lb}_f/\text{ft}^3 = 1278.528 \, \text{lb}_f/\text{ft}^3
\]

\subsubsection*{Part c}

\[
1.61 \, \text{kg/m}^3 = 1.61 \times 0.062428 \, \text{slug/ft}^3 = 0.10051 \, \text{slug/ft}^3
\]

\subsubsection*{Part d}

\[
0.0320 \, \text{N.m/s} = 0.0320 \times 723.569 \, \text{lb}_f.ft/\text{s} = 23.15421 \, \text{lb}_f.ft/\text{s}
\]

\subsubsection*{Part e}

\[
5.67 \, \text{mm/hr} = 5.67 \times 0.00328084 \, \text{in/s} = 0.0186 \, \text{in/s}
\]

\newpage % Starts a new page

\section*{Problem \#4}
An object is observed to weigh 300 N when measured on the surface of the Earth. Determine the mass of the object in kilograms and ascertain its weight in Newtons when placed on the surface of another planet where the acceleration due to gravity is \(4.0 \, \text{ft/s}^2\).

\subsection*{Solution}

Firstly, we establish the known values and constants:
\begin{itemize}
    \item Weight on Earth: \( W = 300 \, \text{N} \)
    \item Acceleration due to gravity on Earth: \( g_{\text{Earth}} = 9.81 \, \text{m/s}^2 \)
    \item Acceleration due to gravity on the other planet: \( g_{\text{Planet}} = 4.0 \, \text{ft/s}^2 \)
    \item Conversion factor for \( \text{ft/s}^2 \) to \( \text{m/s}^2 \): \( 1 \, \text{ft/s}^2 = 0.3048 \, \text{m/s}^2 \)
\end{itemize}

\subsubsection*{Mass of the Object in kg}

We can find the mass of the object using Newton's second law of motion:

\[
m = \frac{W}{g_{\text{Earth}}} = \frac{300 \, \text{N}}{9.81 \, \text{m/s}^2} \approx 30.58 \, \text{kg}
\]

\subsubsection*{Weight on the Other Planet in N}

First, we convert \( g_{\text{Planet}} \) into \( \text{m/s}^2 \):

\[
g_{\text{Planet}} = 4.0 \, \text{ft/s}^2 \times 0.3048 \, \text{m/s}^2/\text{ft/s}^2 = 1.2192 \, \text{m/s}^2
\]
Then we find the weight on the other planet using the mass calculated earlier:

\[
W_{\text{Planet}} = m \times g_{\text{Planet}} = 30.58 \, \text{kg} \times 1.2192 \, \text{m/s}^2 \approx 37.3 \, \text{N}
\]
The mass of the object is approximately \( 30.58 \, \text{kg} \) and its weight on the other planet would be \( 37.3 \, \text{N} \).

\newpage
\section*{Problem \#5}
A tire having a volume of \(3 \, \text{ft}^3\) contains air at a gage pressure of \(26 \, \text{psi}\) and a temperature of \(70 \, \text{F}\). Determine the following:

\begin{enumerate}
    \item The density of the air contained in the tire.
    \item The weight of the air contained in the tire.
\end{enumerate}

\subsection*{Solution}
\subsubsection*{SI Units}

\begin{enumerate}
  \item Absolute Pressure: \( P_{\text{abs}} = 40.7 \, \text{psi} = 280489.5 \, \text{Pa} \)
  \item Temperature: \( T = 70 \, \text{°F} = 294.2611 \, \text{K} \)
  \item Universal Gas Constant: \( R = 287 \, \text{J/(kg$\cdot$K)} \)
  \item Density:
    \[
    \rho = \frac{P_{\text{abs}}}{R \times T} = \frac{280489.5 \, \text{Pa}}{287 \, \text{J/(kg$\cdot$K)} \times 294.2611 \, \text{K}} \approx 3.430847 \, \text{kg/m}^3
    \]
  \item Weight of Air: \( V = 3 \, \text{ft}^3 = 0.084953 \, \text{m}^3 \)
    \[
    W = \rho \times V = 3.430847 \, \text{kg/m}^3 \times 0.084953 \, \text{m}^3 \approx 0.291423 \, \text{kg}
    \]
  \item Conversion to USCS:
    \[
    W_{\text{USCS}} = W_{\text{SI}} \times 2.20462 \, \text{lbm/kg} \approx 0.642372 \, \text{lbm}
    \]
\end{enumerate}

\subsubsection*{USCS Units}

\begin{enumerate}
  \item Absolute Pressure: \( P_{\text{abs}} = 40.7 \, \text{psi} \)
  \item Temperature: \( T = 529.67 \, \text{°R} \)
  \item Universal Gas Constant: \( R = 1716 \, \text{ft$\cdot$lbf/(slug$\cdot$°R)} \)

  \item Conversion Factor:
    \[
    \frac{\text{lbf/in}^2}{\frac{\text{ft$\cdot$lbf}}{\text{slug$\cdot$°R}} \times \text{°R}} = 144 \, \text{slug/ft}^3
    \]
  
  \item Density:
    \[
    \rho = \frac{40.7 \, \text{psi}}{1716 \, \text{ft$\cdot$lbf/(slug$\cdot$°R)} \times 529.67 \, \text{°R}} \times 144 \approx 0.006448 \, \text{slug/ft}^3
    \]
  
  \item Density Conversion:
    \[
    \rho = 0.006448 \, \text{slug/ft}^3 \times 32.174 \, \text{lbm/slug} = 0.207458 \, \text{lbm/ft}^3
    \]
  
  \item Weight of Air:
    \[
    W = \rho \times V = 0.207458 \, \text{lbm/ft}^3 \times 3 \, \text{ft}^3 = 0.622374 \, \text{lbm}
    \]
\end{enumerate}

\newpage
\section*{Problem \#6}
The sled shown in the figure slides along on a thin horizontal layer of water between the ice and the runners. The horizontal force that the water puts on the runners is equal to \(1.2 \, \text{lb}\), when the sled's speed is \(50 \, \text{ft/s}\). The total area of both runners in contact with the water is \(0.08 \, \text{ft}^2\) and the viscosity of the water is \(3.5 \times 10^{-5} \, \text{lb·s/ft}^2\). Determine the thickness of the water layer under the runners, assuming a linear velocity distribution in the water layer. Use USCS units throughout the calculation.

\subsection*{Solution}
To solve this problem, we use the formula for viscous shear force in a fluid layer:

\[
F = \eta \times A \times \frac{\Delta V}{\Delta x}
\]
Given that \( F = 1.2 \, \text{lb}\), \( \eta = 3.5 \times 10^{-5} \, \text{lb·s/ft}^2\), \( A = 0.08 \, \text{ft}^2\), and \( \Delta V = 50 \, \text{ft/s}\), we need to find \( \Delta x \), the thickness of the water layer.
Rearranging the equation to solve for \( \Delta x \):

\[
\Delta x = \eta \times A \times \frac{\Delta V}{F}
\]
Substituting the given values:

\[
\Delta x = 3.5 \times 10^{-5} \, \text{lb·s/ft}^2 \times 0.08 \, \text{ft}^2 \times \frac{50 \, \text{ft/s}}{1.2 \, \text{lb}}
\]
Finally, we get:

\[
\Delta x \approx 1.167 \times 10^{-4} \, \text{ft}
\]
Thus, the thickness of the water layer is approximately \( 1.167 \times 10^{-4} \, \text{ft}\).
\newpage
\section*{Problem \#7}
Some measurements on a blood sample at \(37^\circ \text{C} (98.6^\circ \text{F})\) indicate a shearing stress of \(0.52 \, \text{N/m}^2\) for a corresponding rate of shearing strain of \(200 \, \text{s}^{-1}\). Determine the apparent viscosity of the blood and compare it with the viscosity of water at the same temperature. The viscosity of water at \(37^\circ \text{C}\) is \( \mu_w = 6.9 \times 10^{-4} \, \text{N·s/m}^2\) by linear interpolation.

\subsection*{Solution}
To find the apparent viscosity \( \mu \) of the blood, we use the equation for Newtonian flow:

\[
\tau = \mu \times \gamma
\]
Given \( \tau = 0.52 \, \text{N/m}^2\) and \( \gamma = 200 \, \text{s}^{-1}\), we can rearrange the equation to solve for \( \mu \):

\[
\mu = \frac{\tau}{\gamma}
\]
Substituting the given values:

\[
\mu = \frac{0.52 \, \text{N/m}^2}{200 \, \text{s}^{-1}} = 2.6 \times 10^{-3} \, \text{N·s/m}^2
\]
To compare this with the viscosity of water at the same temperature, \( \mu_w = 6.9 \times 10^{-4} \, \text{N·s/m}^2\):

\[
\frac{\mu}{\mu_w} = \frac{2.6 \times 10^{-3}}{6.9 \times 10^{-4}} \approx 3.77
\]
The blood has an apparent viscosity that is approximately \(3.77\) times higher than that of water at \(37^\circ \text{C}\).
\newpage
\section*{Problem \#8}

During a mountain climbing trip, it is observed that the water used to cook a meal boils at \(90^\circ \text{C}\) rather than the standard \(100^\circ \text{C}\) at sea level. At what altitude are the climbers preparing their meal? See tables B.2 and C.2 for the data needed to solve this problem.

\begin{itemize}
  \item Table B.2: Vapor pressure of water at \(100^\circ \text{C} = 1.013 \times 10^5 \, \text{N/m}^2\)
  \item Table B.2: Vapor pressure of water at \(90^\circ \text{C} = 7.010 \times 10^4 \, \text{N/m}^2\)
  \item Table C.2: Atmospheric pressure at altitude 0 m \(= 1.013 \times 10^5 \, \text{N/m}^2\)
  \item Table C.2: Atmospheric pressure at altitude 3000 m \(= 7.012 \times 10^4 \, \text{N/m}^2\)
  \item Table C.2: Atmospheric pressure at altitude 4000 m \(= 6.166 \times 10^4 \, \text{N/m}^2\)
\end{itemize}

\subsection*{Solutions}

Given that the vapor pressure of water at \(90^\circ \text{C}\) is \(7.010 \times 10^4 \, \text{N/m}^2\), we can equate this to the atmospheric pressure \(P\) at the specific altitude \(h\):

\[
P = 7.010 \times 10^4 \, \text{N/m}^2
\]
From Table C.2, the atmospheric pressure at 3000 m is \(7.012 \times 10^4 \, \text{N/m}^2\), which is very close to \(7.010 \times 10^4 \, \text{N/m}^2\). Therefore, we can approximate that the climbers are at an altitude of 3000 m.
\newpage
\section*{Problem \#9}

When a \(2 \, \text{mm}\) diameter tube is inserted into a liquid in an open tank, the liquid is observed to rise \(10 \, \text{mm}\) above the free surface of the liquid. The contact angle between the liquid and the tube is zero, and the specific weight of the liquid is \(1.2 \times 10^4 \, \text{N/m}^3\). Determine the value of the surface tension for this liquid (N/m).

\subsection*{Solutions}
To find the surface tension \( \sigma \) of the liquid, we use Jurin's Law:

\[
h = \frac{{2 \sigma \cos(\theta)}}{{r \gamma}}
\]
Rearranging the formula to solve for \( \sigma \):

\[
\sigma = \frac{{h \times r \times \gamma}}{2 \cos(\theta)}
\]
Substituting the given values \( h = 0.01 \, \text{m} \), \( r = 0.001 \, \text{m} \), \( \gamma = 1.2 \times 10^4 \, \text{N/m}^3 \), and \( \theta = 0 \):

\[
\sigma = \frac{{0.01 \times 0.001 \times 1.2 \times 10^4}}{2}
\]
\[
= 0.06 \, \text{N/m}
\]
Thus, the surface tension of the liquid is \(0.06 \, \text{N/m}\).
\newpage
\section*{Problem \#10}
The water strider bug shown in the figure is supported on the surface of a pond by surface tension acting along the interface between the water and the bug's legs. Determine the minimum length of this interface needed to support the bug. Assume the bug weighs \(10^{-4} \, \text{N}\) and the surface tension force acts vertically upwards. Assuming \(T = 20^\circ \text{C}\), and \(\sigma = 7.28 \times 10^{-2} \, \text{N/m}\).

\subsection*{Solutions}
To find the minimum length \(L\) of the interface needed to support the bug, we consider the force balance in the vertical direction:

\[
F_{\text{ST}} = F_{\text{g}}
\]
The surface tension force \(F_{\text{ST}}\) is calculated as:

\[
F_{\text{ST}} = \sigma \times L
\]
The gravitational force \(F_{\text{g}}\) is:

\[
F_{\text{g}} = 10^{-4} \, \text{N}
\]
Equating the two forces:

\[
\sigma \times L = 10^{-4} \, \text{N}
\]
Solving for \(L\):

\[
L = \frac{{10^{-4} \, \text{N}}}{{7.28 \times 10^{-2} \, \text{N/m}}}
\]
\[
L \approx 0.001375 \, \text{m} \approx 1.375 \, \text{mm}
\]
Thus, the minimum length of the interface needed to support the bug is approximately \(1.375 \, \text{mm}\).
\newpage
\section*{Problem \#11}

A jet aircraft flies at a speed of 900 mph at an altitude of 30,000 ft, where the temperature is \(60^\circ \text{F}\) and the specific heat ratio is \(k=1.4\). Determine the Mach number, \( \text{Ma} \) (the ratio of the speed of the aircraft, \( V \), to that of the speed of sound, \( c \)) at the specified altitude. The speed of sound is given as follows: \( c = \sqrt{kRT} \).

\subsection*{Solutions}

To find the Mach number at the specified altitude, first we convert the speed of the aircraft \( V \) to ft/s:

\[
V = 900 \times 1.467 = 1320.3 \, \text{ft/s}
\]
Next, we calculate the speed of sound \( c \) using the updated value for \( R \) of \( 1716 \, \text{ft} \cdot \text{lbf/(slug} \cdot \text{R)} \):

\[
c = \sqrt{1.4 \times 1716 \times 519.67} \approx 1116.5 \, \text{ft/s}
\]
Finally, the Mach number \( \text{Ma} \) is calculated as:

\[
\text{Ma} = \frac{V}{c} = \frac{1320.3}{1116.5} \approx 1.183
\]
Thus, the Mach number of the aircraft at the specified altitude is approximately \( \text{Ma} = 1.183 \).
\newpage
\section*{Problem \#12}
A regulation basketball is initially flat and then is inflated to a pressure of approximately \(24 \, \text{lb/in}^2\) absolute. Considering the air temperature to be constant at \(70^\circ \text{F}\), find the mass of the air required to inflate the basketball. The basketball's inside radius is \(4.67 \, \text{in}\).

\subsection*{Solutions}
To determine the mass of the air needed, we use the ideal gas law:

\[
PV = mRT
\]
We're given the pressure \( P = 24 \, \text{lb/in}^2 \), the specific gas constant \( R = 1716 \, \text{ft} \cdot \text{lbf/(slug} \cdot \text{R)} \), and the temperature \( T = 70 + 459.67 = 529.67 \, \text{°R} \).
\newline
\newline
First, convert the pressure and volume into consistent units:

\[
P = 24 \, \text{lb/in}^2 \times \left( \frac{144 \, \text{in}^2}{1 \, \text{ft}^2} \right) = 3456 \, \text{lb/ft}^2
\]

\[
V = \frac{4}{3} \pi (4.67 \, \text{in})^3 \times \left( \frac{1 \, \text{ft}^3}{1728 \, \text{in}^3} \right) = 0.248 \, \text{ft}^3
\]
Finally, solving for the mass \( m \) gives:

\[
m = \frac{PV}{RT} = \frac{3456 \times 0.248}{1716 \times 529.67} \approx 0.00109 \, \text{slug}
\]
To convert the mass to pound-mass (lbm), we use the following relationship:

\[
m = 0.00109 \, \text{slug} \times 32.174 \, \text{lbm/slug} \approx 0.0351 \, \text{lbm}
\]
Thus, the mass of the air required to inflate the basketball is approximately \(0.0351 \, \text{lbm}\).


\end{document}
